\documentclass[11pt,a4paper]{article}

% ============================================================
% Layout and typography
% ============================================================
\usepackage[T1]{fontenc}
\usepackage[utf8]{inputenc}
\usepackage{lmodern}

\usepackage[
  a4paper,
  top=1.9cm,
  bottom=2.0cm,
  left=2.1cm,
  right=2.1cm
]{geometry}

\usepackage{microtype}
\usepackage{setspace}

\setstretch{1.03}
\setlength{\parindent}{0pt}
\setlength{\parskip}{0.32em}
\setlength{\emergencystretch}{2em}

% ============================================================
% Mathematics, tables and figures
% ============================================================

\usepackage{amsmath,amssymb}
\usepackage{booktabs}
\usepackage{tabularx}
\usepackage{array}
\usepackage{graphicx}
\usepackage{float}
\usepackage{placeins}
\usepackage{siunitx}
\usepackage{enumitem}
\usepackage{listings}

\sisetup{
  round-mode=places,
  round-precision=4,
  detect-all,
  group-digits=false
}

\setlist[itemize]{leftmargin=1.35em,itemsep=0.15em,topsep=0.2em}
\setlist[enumerate]{leftmargin=1.55em,itemsep=0.15em,topsep=0.2em}

% ============================================================
% Colors, captions and hyperlinks
% ============================================================

\usepackage{xcolor}

\definecolor{linkblue}{RGB}{28,72,122}

\usepackage[
font=small,
labelfont=bf,
justification=centering,
singlelinecheck=false,
skip=6pt
]{caption}

\usepackage[
colorlinks=true,
linkcolor=linkblue,
citecolor=linkblue,
urlcolor=linkblue,
pdfauthor={Momen Ismail},
pdftitle={Comparing Machine Learning Models for Monthly Stock Return Prediction}
]{hyperref}

\usepackage[nameinlink,noabbrev]{cleveref}

% ============================================================
% Bibliography
% ============================================================

\usepackage{csquotes}

\usepackage[
backend=biber,
style=authoryear,
maxcitenames=2,
giveninits=true,
uniquename=init,
doi=false,
url=false,
isbn=false
]{biblatex}

\addbibresource{references.bib}

% ============================================================
% Header
% ============================================================

\usepackage{fancyhdr}

\pagestyle{fancy}
\fancyhf{}
\fancyhead[L]{Machine Learning for Finance}
\fancyhead[R]{Momen Ismail}
\fancyfoot[C]{\thepage}
\renewcommand{\headrulewidth}{0.4pt}
\setlength{\headheight}{14pt}

% ============================================================
% Section style
% ============================================================

\usepackage{titlesec}

\titleformat{\section}
{\large\bfseries}
{\thesection}
{0.65em}
{}

\titleformat{\subsection}
{\normalsize\bfseries}
{\thesubsection}
{0.60em}
{}

\titlespacing*{\section}{0pt}{1.0em}{0.25em}
\titlespacing*{\subsection}{0pt}{0.70em}{0.12em}

\setcounter{secnumdepth}{2}

% ============================================================
% Project paths
% ============================================================

\newcommand{\FigurePath}{../../figures}
\newcommand{\papertitle}{Comparing Machine Learning Models for Monthly Stock Return Prediction}

% ============================================================
% Document
% ============================================================

\begin{document}

% ============================================================
% Title page (excluded from the 15-page limit)
% ============================================================

\begin{titlepage}
\thispagestyle{empty}
\centering

\vspace*{2.2cm}

{\LARGE\bfseries
\papertitle
\par}

\vspace{1.2cm}

{\Large Course Project Report\par}

\vspace{1.4cm}

{\large
Momen Ismail\par
}

\vspace{0.35cm}

{\large
Matriculation Number: 50307564\par
}

\vspace{0.35cm}

{\large
Institutional Email:\\[0.15cm]
\texttt{\href{mailto:s73misma@uni-bonn.de}{s73misma@uni-bonn.de}}
\par
}

\vspace{1.2cm}

{\large
Machine Learning for Finance\par
University of Bonn\par
Instructor: Prof. Ivan Gufler
\par
}

\vfill

\textbf{GitHub Repository}

\vspace{0.25cm}

\url{https://github.com/Momen-Ismail/stock-return-ml-project}

\vfill

27 July 2026

\end{titlepage}

\clearpage
\pagenumbering{arabic}

\begin{abstract}
This project compares six models for predicting next-month stock excess
returns: a historical mean benchmark, OLS-3, Partial Least Squares, Elastic
Net, Random Forest, and Gradient Boosting. I use a monthly panel with 211,059
stock-month observations, 656 tickers, and 484 predictors from February 1990
to December 2025. The model choices are made using chronological training and
validation periods, and the final specifications are evaluated once on the
2020--2025 test sample. Random Forest performs best and reduces monthly squared
forecast error by about 1.14\% relative to the historical mean. OLS-3 provides
a very small improvement, Elastic Net reduces all coefficients to zero, and
PLS and Gradient Boosting perform worse than the benchmark. The main conclusion
is that machine learning can provide a modest improvement, but the result is
not stable across years and does not by itself imply a profitable trading
strategy.
\end{abstract}

\noindent\textbf{Keywords:} stock return prediction, machine learning,
out-of-sample evaluation, Random Forest, asset pricing

% ============================================================
% 1. Introduction
% ============================================================
\section{Introduction}
\label{sec:introduction}

Predicting stock returns is a difficult problem. Monthly returns are noisy,
market conditions change, and relationships that appear useful in one period
may disappear in another. At the same time, investors and researchers have a
large amount of information available, including price history, accounting
variables, market conditions, and macroeconomic indicators. This creates a
natural setting for comparing traditional statistical models with machine
learning methods.

Machine learning is useful in empirical asset pricing for three main reasons.
First, expected-return estimation is a prediction problem: information known at
time \(t\) is used to predict returns at time \(t+1\). Second, the number of
possible predictors is large and many variables are correlated. Third, the
relationship between characteristics and future returns may be nonlinear and
may depend on interactions between firm-level and aggregate variables. These
points motivate regularization, dimension reduction, and tree-based methods
\parencite{gu2020empirical}.

The research question of this project is:

\begin{quote}
\emph{Can machine learning models predict next-month stock excess returns more
accurately than a simple historical-mean benchmark?}
\end{quote}

I compare six models under the same data and sample design. The historical mean
is the benchmark. OLS-3 is a simple linear model using size, book-to-market, and
12-month momentum. PLS and Elastic Net use all 484 predictors but control the
high dimensionality in different ways. Random Forest and Gradient Boosting
allow nonlinear relationships and interactions.

The comparison is designed to be realistic for financial prediction. I do not
use random train-test splits. Hyperparameters are selected using expanding
windows inside the training period, fixed and optimized versions are compared
on a later validation period, and the final models are evaluated once on the
2020--2025 test period. This follows the main lesson from the forecasting
literature: a model should be judged on genuinely future observations and
against a clear benchmark \parencite{welch2008comprehensive}.

The project is mainly a comparison exercise. I do not propose a new algorithm,
and I do not claim that the best model creates a profitable trading strategy.
The goal is to see which model generalizes best, where the models fail, and
whether greater complexity is actually useful in this data set.

The rest of the paper is organized as follows. \Cref{sec:data} describes the
data and preprocessing. \Cref{sec:methodology} explains the models and the
comparison design. \Cref{sec:results} presents the validation and test results.
\Cref{sec:critical-reflection} discusses the main limitations and what I would
do differently. \Cref{sec:conclusion} concludes.

% ============================================================
% 2. Data
% ============================================================
\section{Data}
\label{sec:data}

\subsection{Sources and observation unit}

The observation unit is a stock-month. The stock universe is based on a locked
S\&P 500-style list that contains current and historical tickers available in
the project inputs. This is not a perfect point-in-time index-constituent data
set, but the same locked universe is used throughout the pipeline.

Daily stock prices and trading volume come from Yahoo Finance. Daily S\&P 500
and VIX data are used for market and volatility variables. Monthly
Fama--French data provide the risk-free rate. Annual Compustat data provide
accounting variables and industry classifications. Eight aggregate predictor
states are taken from the Welch--Goyal data set.

Daily histories begin in 1987 so that longer-window variables can be calculated
before the main sample starts. The final model-ready panel covers February 1990
to December 2025.

\subsection{Prediction target}

For stock \(i\) in month \(t\), the target is the next-month excess return:

\begin{equation}
 y_{i,t+1}=r_{i,t+1}-RF_{t+1},
 \label{eq:target}
\end{equation}

where \(r_{i,t+1}\) is the stock return in the next calendar month and
\(RF_{t+1}\) is the corresponding risk-free rate.

The target is created only when the next observation for the same ticker is the
immediately following calendar month. This avoids treating a return over a gap
of several months as a one-month return. The target itself is kept in its raw
return units and is not winsorized.

\subsection{Timing and leakage controls}

The main timing rule is simple: a row dated month \(t\) may only use information
available by the end of month \(t\).

Stock-price characteristics use price history through month \(t\). Compustat
requires an additional timing rule because the fiscal year-end date is not the
same as the public reporting date. I therefore treat annual accounting data as
available six months after fiscal year-end and use the latest report available
for each stock-month. Very old accounting observations are not carried forward
indefinitely.

The Welch--Goyal variables are lagged by one month. Therefore, a stock-month in
month \(t\) uses the aggregate state from month \(t-1\). Together with the
exact next-month target rule, these choices reduce look-ahead bias.

\subsection{Preprocessing}

Missing values in the 45 continuous firm characteristics are replaced with the
cross-sectional median of the same month. This uses information from other
firms in that month, but not from future months. The continuous characteristics
are then winsorized within each month at the 1st and 99th percentiles. This
limits the influence of extreme predictor values while keeping the target
unchanged.

The three binary characteristics remain binary. Market and VIX variables,
Welch--Goyal variables, and SIC2 indicators are not cross-sectionally
winsorized. The 360 interaction variables are formed from the processed
continuous characteristics and the eight aggregate states.

\subsection{Predictors and final sample}

The final data contain 484 predictors:

\begin{itemize}
  \item 45 continuous stock and accounting characteristics;
  \item 3 binary stock characteristics;
  \item 4 market and VIX variables;
  \item 8 lagged Welch--Goyal variables;
  \item 64 SIC2 industry indicators;
  \item 360 characteristic--macro interactions.
\end{itemize}

The continuous characteristics cover size, valuation, profitability,
investment, leverage, liquidity, momentum, reversal, volatility, beta, and
trading activity. The interactions allow the effect of a firm characteristic
to depend on the aggregate economic state.

The final panel contains 211,059 observations and 656 tickers. It has no missing
predictor values, infinite values, or duplicated ticker-month observations.

\begin{table}[H]
\centering
\caption{Chronological sample design}
\label{tab:sample-splits}
\small
\begin{tabular}{lrrrrl}
\toprule
Sample & Observations & Months & Tickers & Start & End \\
\midrule
Training    & 131,061 & 299 & 572 & 1990-02 & 2014-12 \\
Validation  &  35,386 &  60 & 611 & 2015-01 & 2019-12 \\
Development & 166,447 & 359 & 613 & 1990-02 & 2019-12 \\
Test        &  44,612 &  72 & 641 & 2020-01 & 2025-12 \\
\bottomrule
\end{tabular}
\end{table}

The training sample is used for tuning. The validation sample is used to choose
between fixed and optimized versions. The development sample combines training
and validation for the final fit. The test sample is used only for the final
comparison.

% ============================================================
% 3. Methodology
% ============================================================

\section{Methodology}
\label{sec:methodology}

\subsection{Prediction framework}

The purpose of the models is to use information known at the end of month
\(t\) to predict the excess return realized in month \(t+1\). I write the
forecasting problem as

\begin{equation}
 r^e_{i,t+1}
 =
 \mathbb{E}_t
 \left[
 r^e_{i,t+1}
 \right]
 +
 \varepsilon_{i,t+1},
 \qquad
 \mathbb{E}_t
 \left[
 r^e_{i,t+1}
 \right]
 =
 g(\mathbf{x}_{i,t}),
 \label{eq:framework}
\end{equation}

where \(r^e_{i,t+1}\) is the next-month excess return of stock \(i\),
\(\mathbf{x}_{i,t}\) is the information available at month \(t\), and
\(\varepsilon_{i,t+1}\) is the part of the return that is not predicted by the
model.

The models differ in how they estimate the function \(g(\cdot)\). Some impose a
simple linear relationship, while others allow regularization, dimension
reduction, nonlinearities, and interactions. I evaluate all models using the
same target, chronological samples, and performance measures so that their
results are directly comparable.

\subsection{Model choice and rationale}

I chose six models that represent different levels and types of model
complexity. The purpose is not only to find the best-performing model, but also
to test whether additional complexity is actually useful for this prediction
problem.

\begin{table}[H]
\centering
\caption{Models included in the comparison}
\label{tab:model-rationale}
\small
\begin{tabularx}
{\textwidth}
{>{\raggedright\arraybackslash}p{2.8cm}X}
\toprule
Model & Reason for inclusion \\
\midrule

Historical Mean &
A simple constant benchmark. It shows whether the other models provide any
improvement over the average return observed in the available estimation
sample. \\

OLS-3 &
A small and interpretable linear model using size, book-to-market, and
12-month momentum. It provides a comparison between a simple economic model
and the high-dimensional methods. \\

Partial Least Squares &
A dimension-reduction model that replaces the 484 correlated predictors with a
small number of supervised components. \\

Elastic Net &
A regularized linear model that can shrink unstable coefficients and set some
of them exactly to zero. \\

Random Forest &
A tree ensemble that can capture nonlinear relationships and interactions while
reducing the instability of an individual tree through averaging. \\

Gradient Boosting &
A second tree ensemble that fits shallow trees sequentially, with each new tree
trying to correct errors made by the previous trees. \\

\bottomrule
\end{tabularx}
\end{table}

\subsubsection*{Historical Mean}

The historical mean is the main benchmark. It ignores all stock
characteristics and assigns the same prediction to every stock. During each
training or validation exercise, the mean is calculated using only the
permitted estimation sample. For the final test, it is calculated using the
complete February 1990--December 2019 development sample.

This is an important benchmark because a machine-learning model is only useful
for this project if it improves on a simple constant forecast.

\subsubsection*{OLS-3}

OLS-3 uses three predictors:

\begin{itemize}
  \item firm size;
  \item book-to-market;
  \item 12-month momentum.
\end{itemize}

I chose these variables because they are well-known and economically
interpretable stock characteristics
\parencite{fama1992cross,jegadeesh1993returns}. The model provides a useful
middle point between the historical mean and the 484-predictor models.

It also allows me to examine whether a small set of established predictors can
perform as well as, or better than, much more complex methods.

\subsubsection*{Partial Least Squares}

Partial Least Squares (PLS) is included because the predictor set is
high-dimensional and strongly correlated. PLS constructs a small number of
latent components using information from both the predictor matrix and the
target. It therefore differs from an unsupervised method such as ordinary
Principal Component Analysis, which constructs components without using the
target variable \parencite{wold1984collinearity}.

The main hyperparameter is the number of PLS components. Too few components may
remove useful information, while too many components may reproduce noise from
the original predictor set.

\subsubsection*{Elastic Net}

Elastic Net is included as a high-dimensional linear model. It combines two
types of regularization:

\begin{itemize}
  \item the \(L_1\) penalty, which can set coefficients to zero;
  \item the \(L_2\) penalty, which stabilizes coefficients when predictors are
        strongly correlated.
\end{itemize}

The parameter \(\alpha\) controls the overall strength of regularization, while
the \(l_1\)-ratio determines the balance between the two penalties
\parencite{tibshirani1996regression,zou2005regularization}.

This model is useful because it tests whether the predictive information can be
represented through a sparse or heavily regularized linear relationship.

\subsubsection*{Random Forest}

Random Forest fits many regression trees to bootstrap samples and averages
their predictions \parencite{breiman2001random}. Each split considers only a
random subset of predictors.

The random sampling of both observations and predictors reduces the correlation
between the trees. Averaging then reduces the variance of an individual tree.
The model can still represent nonlinearities and interactions without requiring
these relationships to be specified in advance.

The Random Forest specification also uses out-of-bag observations. I use the
out-of-bag result only as an additional development-sample diagnostic. It is not
treated as final test evidence and is not used instead of the chronological
validation and test samples.

\subsubsection*{Gradient Boosting}

Gradient Boosting also combines regression trees, but it does so sequentially.
Each new tree is fitted to improve the errors remaining from the existing
ensemble.

The main decisions are:

\begin{itemize}
  \item the number of trees;
  \item the learning rate;
  \item the maximum depth of each tree.
\end{itemize}

A small learning rate makes each tree contribute only a limited correction.
Shallow trees restrict the complexity of each individual learner. The number of
trees determines how many sequential corrections the model is allowed to make.

\subsection{Model-specific scaling and implementation choices}

The continuous firm characteristics were already median-imputed and winsorized
during data construction. For model estimation, I make an additional scaling
decision.

The predictors used by OLS-3, PLS, and Elastic Net are standardized using the
mean and standard deviation of the relevant estimation sample. The same fitted
transformation is then applied to the corresponding validation or test sample.

Standardization is important for PLS and Elastic Net because their components
and penalties depend on the scale of the predictors. I also standardize OLS-3
for consistency, although standardization does not change the fitted values of
an unregularized OLS regression with an intercept.

Random Forest and Gradient Boosting are not standardized. Tree splits depend on
the ordering of predictor values rather than their measurement units, so
rescaling the variables is unnecessary.

Both tree models use a fixed random seed of 42 so that the results can be
reproduced.

\subsection{Tuning design}

I do not use random cross-validation. A random split would mix earlier and
later market observations and could allow future market conditions to influence
model selection.

Instead, hyperparameters are selected using ten annual expanding-window folds
inside the February 1990--December 2014 training sample:

\begin{itemize}
  \item the first fold trains through December 2004 and validates on 2005;
  \item the training window then expands by one year for each new fold;
  \item the final fold trains through December 2013 and validates on 2014.
\end{itemize}

For each candidate specification, I calculate monthly MSE in every validation
year and then average the ten annual results:

\begin{equation}
 \operatorname{CV\text{-}MSE}(\lambda)
 =
 \frac{1}{10}
 \sum_{k=1}^{10}
 \operatorname{MSE}_{\mathrm{month},k}(\lambda),
 \label{eq:cv-mse}
\end{equation}

where \(\lambda\) represents the hyperparameter combination.

The tuning stage examines:

\begin{itemize}
  \item the number of components for PLS;
  \item \(\alpha\) and the \(l_1\)-ratio for Elastic Net;
  \item the number of trees for Random Forest;
  \item the number of trees, learning rate, and maximum depth for Gradient
        Boosting.
\end{itemize}

I use monthly MSE as the tuning criterion because the number of stocks differs
across months. Averaging errors within month first prevents months with more
observations from receiving greater weight.

\subsection{Fixed and optimized specifications}

Before tuning, I estimate fixed reference versions of the models. These provide
a direct test of whether optimization actually improves performance.

The main fixed specifications are:

\begin{itemize}
  \item PLS with 20 components;
  \item Elastic Net with \(\alpha=0.0001\) and \(l_1\)-ratio \(=0.50\);
  \item Random Forest with 100 trees;
  \item Gradient Boosting with 100 trees, learning rate \(0.01\), and maximum
        depth 2.
\end{itemize}

The expanding-window training exercise produces optimized candidates. However,
I do not automatically send these candidates to the final test. I first compare
each optimized specification with its fixed version using the later
January 2015--December 2019 validation sample.

This additional validation step is important because a specification that
performs well in earlier training folds may not remain useful in a later market
period.

\subsection{Validation-based final model selection}

The validation results are shown in
\Cref{tab:validation-choice}.

\begin{table}[H]
\centering
\caption{Validation choice between fixed and optimized specifications}
\label{tab:validation-choice}
\small
\setlength{\tabcolsep}{4.5pt}

\begin{tabularx}
{\textwidth}
{>{\raggedright\arraybackslash}p{2.5cm}
 S[table-format=1.6]
 S[table-format=1.6]
 >{\raggedright\arraybackslash}X}
\toprule
{Model family}
& {Fixed monthly MSE}
& {Optimized monthly MSE}
& {Final validation choice} \\
\midrule

PLS
& 0.007710
& 0.007627
& Optimized version with 2 components \\

Elastic Net
& 0.007750
& 0.007449
& Optimized version with \(\alpha=0.015\) and \(l_1\)-ratio \(=0.85\) \\

Random Forest
& 0.007516
& 0.007534
& Fixed version with 100 trees \\

Gradient Boosting
& 0.007450
& 0.007588
& Fixed version with 100 trees, learning rate \(0.01\), and depth 2 \\

\bottomrule
\end{tabularx}
\end{table}

The results lead to four different decisions.

First, the optimized PLS model is retained because two components produce a
lower validation MSE than the fixed 20-component model. This suggests that a
small number of supervised components generalizes better than a larger
component representation.

Second, the optimized Elastic Net is retained because it improves on the fixed
Elastic Net in validation. The selected penalty is strong enough to shrink all
484 coefficients to zero. I do not treat this as a coding failure. It means
that, under the selected chronological validation design, the model preferred a
constant prediction over unstable linear predictor effects.

Third, expanding-window tuning favored a larger Random Forest candidate with
300 trees, but the fixed 100-tree model produces a lower MSE in the separate
2015--2019 validation period. I therefore retain the fixed 100-tree version.

Fourth, the more complex 300-tree Gradient Boosting candidate also performs
worse than the fixed 100-tree model in the later validation period. I therefore
retain the fixed version.

These decisions show why I use both tuning folds and a separate validation
sample. Optimization inside the training period does not automatically imply
better performance in a genuinely later period.

The final Random Forest specification is:

\begin{itemize}
  \item 100 trees;
  \item square-root feature sampling at each split;
  \item minimum leaf size of 20;
  \item bootstrap sampling;
  \item out-of-bag diagnostics;
  \item random seed 42.
\end{itemize}

The final Gradient Boosting specification is:

\begin{itemize}
  \item 100 trees;
  \item learning rate \(0.01\);
  \item maximum tree depth 2;
  \item random seed 42.
\end{itemize}

All final choices are made using the training and validation periods. The
2020--2025 test results are not used to choose these specifications.

\subsection{Final estimation and test procedure}

After the final specifications are selected, I combine the training and
validation samples into the February 1990--December 2019 development sample.

The final procedure is:

\begin{enumerate}
  \item select hyperparameters using the 1990--2014 training sample;
  \item compare fixed and optimized models using the 2015--2019 validation
        sample;
  \item lock the final specifications;
  \item refit each selected model on the complete 1990--2019 development
        sample;
  \item evaluate each fitted model once on the January 2020--December 2025
        test sample.
\end{enumerate}

The final test contains 44,612 stock-month observations over 72 months. After
the test results are opened, I do not change the predictors, hyperparameters,
sample dates, or model specifications.

\subsection{Evaluation metrics}

The primary evaluation measure is monthly mean squared error:

\begin{equation}
 \operatorname{MSE}_{\mathrm{month}}
 =
 \frac{1}{T}
 \sum_{t=1}^{T}
 \left[
 \frac{1}{N_t}
 \sum_{i=1}^{N_t}
 \left(
 y_{i,t+1}
 -
 \widehat{y}_{i,t+1}
 \right)^2
 \right].
 \label{eq:monthly-mse}
\end{equation}

The squared errors are first averaged across stocks within each month and then
across months. This gives every month the same weight, even when the number of
stocks differs.

The main benchmark-relative measure is monthly out-of-sample \(R^2\):

\begin{equation}
 R^2_{\mathrm{OOS,month}}
 =
 1
 -
 \frac{
   \operatorname{MSE}_{\mathrm{month,model}}
 }{
   \operatorname{MSE}_{\mathrm{month,HM}}
 }.
 \label{eq:monthly-oos-r2}
\end{equation}

A positive value means that the model reduces average monthly squared forecast
error relative to the historical mean. A negative value means that it performs
worse than the benchmark.

The benchmark in this project is the historical mean estimated from the
permitted estimation sample. It is not a zero-return forecast. Therefore, the
out-of-sample \(R^2\) measures improvement relative to an estimated constant
expected-return forecast.

I also report:

\begin{itemize}
  \item monthly RMSE;
  \item pooled RMSE;
  \item pooled MAE;
  \item pooled out-of-sample \(R^2\);
  \item prediction-target correlation.
\end{itemize}

The monthly measures give equal weight to months, while the pooled measures
give equal weight to individual stock-month observations. I treat the monthly
measure as the main criterion and the pooled measures as supporting evidence.

Correlation is undefined for constant predictions. This applies to the
historical mean and to the final Elastic Net model, which produces an
effectively constant prediction.

\subsection{Scope of the comparison}

The model ranking is descriptive. I do not conduct a formal statistical test
showing that the leading model is significantly better than every competing
model.

I also do not construct investment portfolios or calculate turnover,
transaction costs, Sharpe ratios, or cumulative trading returns. The
methodology therefore compares statistical forecast accuracy, not investment
profitability.

Finally, I do not create an additional ensemble that combines the predictions
of all six models. Random Forest and Gradient Boosting are themselves ensemble
methods, but the final comparison evaluates each of the six approaches
separately.

% ============================================================
% 4. Results
% ============================================================

\section{Results}
\label{sec:results}

\subsection{Overall final-test performance}

After completing the tuning and validation stages, I locked the final model
specifications and refitted them on the complete February 1990--December 2019
development sample. I then evaluated them once on the January 2020--December
2025 test sample.

\Cref{tab:test-results} reports the final results. The models are ranked by
monthly MSE, which is the main evaluation criterion. The two out-of-sample
\(R^2\) columns are reported in percent to make the relatively small
differences easier to read.

\begin{table}[H]
\centering
\caption{Final test performance, 2020--2025}
\label{tab:test-results}

\small
\setlength{\tabcolsep}{3.8pt}

\begin{tabular}{clrrrrr}
\toprule
Rank
& Model
& Monthly MSE
& Monthly RMSE
& Monthly OOS \(R^2\) (\%)
& Pooled OOS \(R^2\) (\%)
& Correlation \\
\midrule

1 & Random Forest
  & 0.022112 & 0.148700 &  1.1427 &  1.1071 &  0.1138 \\

2 & OLS-3
  & 0.022360 & 0.149534 &  0.0301 &  0.0296 &  0.0313 \\

3 & Elastic Net
  & 0.022367 & 0.149556 &  0.0000 &  0.0000 &  -- \\

4 & Historical Mean
  & 0.022367 & 0.149556 &  0.0000 &  0.0000 &  -- \\

5 & PLS
  & 0.022509 & 0.150032 & -0.6363 & -0.6579 &  0.0270 \\

6 & Gradient Boosting
  & 0.022573 & 0.150242 & -0.9192 & -0.9203 & -0.0481 \\

\bottomrule
\end{tabular}

\vspace{0.3em}

\begin{minipage}{0.94\textwidth}
\footnotesize
\textit{Notes:}
Models are ranked by monthly MSE. The reported out-of-sample \(R^2\) values
have been multiplied by 100 and are therefore shown as percentages. Monthly
OOS \(R^2\) compares monthly MSE with the historical-mean benchmark. Pooled
OOS \(R^2\) uses all stock-month squared errors. Correlation is undefined for
constant predictions.
\end{minipage}
\end{table}

Random Forest ranks first according to the prespecified monthly MSE criterion.
Its monthly out-of-sample \(R^2\) is \(1.1427\%\). This means that the model
reduces average monthly squared prediction error by approximately \(1.14\%\)
relative to the historical mean. It does not mean that the model predicts a
return of \(1.14\%\).

The pooled out-of-sample \(R^2\) is similar at \(1.1071\%\). The monthly measure
gives every month equal weight, while the pooled measure gives every
stock-month observation equal weight. The similar values show that the result
is not driven only by months containing more stocks.

OLS-3 ranks second, but its improvement is very small. Its monthly
out-of-sample \(R^2\) is only \(0.0301\%\). Nevertheless, this is an important
comparison because OLS-3 uses only size, book-to-market and 12-month momentum.
A simple three-variable model therefore remains competitive with the much
larger models.

Elastic Net is effectively tied with the historical mean. The selected penalty
shrinks all 484 coefficients to zero, so the model produces a constant
prediction. I interpret this as a valid negative result rather than a coding
error. Under the selected penalty and chronological validation design, Elastic
Net did not find stable linear effects that improved on the benchmark.

PLS and Gradient Boosting both have negative out-of-sample \(R^2\). Their
predictions therefore produce greater squared error than the historical mean.
Gradient Boosting also has a negative prediction-target correlation. These
results show that additional model complexity does not automatically improve
out-of-sample performance.

\begin{figure}[H]
\centering
\includegraphics[
  width=0.86\textwidth
]{\FigurePath/test_monthly_oos_r2_by_model.png}

\caption{
Monthly out-of-sample \(R^2\) relative to the historical-mean benchmark in the
2020--2025 test sample. The graph displays the \(R^2\) values in decimal form,
while \Cref{tab:test-results} reports them in percent. Random Forest's value of
approximately \(0.011427\) corresponds to an error reduction of approximately
\(1.1427\%\).
}
\label{fig:test-oos-r2}
\end{figure}

\Cref{fig:test-oos-r2} makes the benchmark comparison visible. Random Forest
has the clearest positive result. OLS-3 is only slightly above zero, Elastic
Net overlaps with the historical mean, and PLS and Gradient Boosting are below
zero.

The ranking is descriptive. I do not estimate confidence intervals for the
differences between models and I do not conduct a multiple-model forecast
comparison test. I can therefore say that Random Forest ranks first according
to monthly MSE, but I cannot claim that it is statistically significantly
better than every alternative.

\subsection{Performance across test years}

The complete test-period result does not mean that Random Forest performs well
in every year. \Cref{fig:yearly-oos-r2} reports the monthly out-of-sample
\(R^2\) separately for each of the six test years.

\begin{figure}[H]
\centering
\includegraphics[
  width=0.92\textwidth
]{\FigurePath/yearly_monthly_oos_r2_by_model.png}

\caption{
Year-specific monthly out-of-sample \(R^2\) for 2020--2025. Random Forest
ranks first over the complete test period but has a positive annual value in
only one of the six years. The historical mean and the near-constant Elastic
Net predictions overlap close to zero.
}
\label{fig:yearly-oos-r2}
\end{figure}

Random Forest has a positive annual out-of-sample \(R^2\) only in 2020. Its
result in that year is sufficiently strong to offset several negative annual
values and produce a positive result over the complete test period.

OLS-3 has a much smaller full-period improvement, but its annual
out-of-sample \(R^2\) is positive in four of the six test years. PLS and
Gradient Boosting are positive in only one year each.

This comparison changes how I interpret the overall ranking. Random Forest has
the lowest average forecast error over the complete period, but its advantage
is concentrated in one year and is not stable across all market environments.
OLS-3 produces a much smaller average gain, but its yearly performance is more
frequently above the benchmark.

The yearly results therefore support a cautious conclusion. Random Forest
performs best according to the main full-period measure, but the evidence does
not show that it will dominate the simpler models in every future period.

\subsection{What the Random Forest uses}

Because Random Forest is the highest-ranked model, I also examine where its
fitted prediction variation comes from. Individual predictor names are often
technical, especially for the interaction variables, so I first summarize
importance by broad predictor group.

\begin{figure}[H]
\centering
\includegraphics[
  width=0.82\textwidth
]{\FigurePath/random_forest_predictor_group_importance.png}

\caption{
Random Forest impurity importance aggregated by predictor group in the final
development-sample fit. The values describe the fitted model and should not be
interpreted as causal effects or estimated risk premia.
}
\label{fig:rf-group-importance}
\end{figure}

Characteristic--macro interactions account for approximately \(73.7\%\) of
the total Random Forest impurity importance. The Welch--Goyal variables account
for approximately \(11.7\%\), market and VIX variables for \(7.8\%\), and the
remaining stock characteristics for approximately \(6.8\%\). SIC2 industry
indicators contribute almost no measured importance.

This pattern suggests that Random Forest obtains much of its fitted prediction
variation from relationships that change with aggregate economic conditions.
In other words, the model appears to use combinations of firm characteristics
and macroeconomic states rather than relying on one single characteristic.

However, this conclusion has to be treated carefully. The interaction group
contains 360 predictors, while the other groups contain far fewer variables.
Its high total importance therefore does not mean that every interaction is
individually important.

Impurity importance also has general limitations. Correlated variables can
divide importance among themselves, and the values can depend on the number of
possible split points. The figure describes how the fitted Random Forest makes
predictions; it does not identify causal effects.

\subsection{Comparison with the machine-learning asset-pricing literature}

The general motivation of this project follows the machine-learning
asset-pricing literature. Flexible models may be useful because they can handle
many predictors and allow nonlinear relationships and interactions
\parencite{gu2020empirical}.

The results of this project are mixed. Random Forest improves on the
historical-mean benchmark, while Gradient Boosting and PLS do not. OLS-3 also
remains competitive despite using only three predictors. The results therefore
support the idea that flexible models can help, but they also show that greater
complexity is not always rewarded outside the estimation sample.

The numerical results should not be directly compared with percentages reported
in another study. The stock universe, sample period, predictors, preprocessing,
and benchmark can differ. In particular, the out-of-sample \(R^2\) in this
project uses the estimated historical mean as its benchmark, rather than a
zero-return benchmark.

\subsection{Financial interpretation}

The project compares statistical forecast accuracy. It does not evaluate a
trading strategy.

A higher out-of-sample \(R^2\) does not automatically imply a higher Sharpe
ratio or a profitable long-short portfolio. A model may reduce average squared
prediction error without producing a useful ranking after trading costs and
portfolio constraints are considered.

The current analysis does not include:

\begin{itemize}
  \item portfolio formation;
  \item transaction costs;
  \item portfolio turnover;
  \item liquidity or market-impact costs;
  \item short-selling restrictions;
  \item risk-adjusted portfolio returns;
  \item Sharpe ratios or cumulative trading returns.
\end{itemize}

The financial conclusion is therefore limited. Random Forest has the lowest
forecast error in the 2020--2025 test period, but the project does not show that
trading on its predictions would be profitable.

A portfolio analysis would be a separate extension. It would require a
predefined sorting or weighting rule and an evaluation of both gross and net
returns. It should not be designed after examining the same test results,
because that would use the test sample for a second model-selection exercise.

% ============================================================
% 5. Critical Reflection
% ============================================================

\section{Critical Reflection}
\label{sec:critical-reflection}

This project involved more than estimating and comparing six models. A large
part of the work was identifying problems in the data pipeline, checking whether
the timing rules were realistic, and deciding which results could be treated as
final. In this section, I discuss the main problems I encountered, how I dealt
with them, and which limitations remain despite the corrections.

\subsection{Problems encountered and how I addressed them}

\begin{table}[H]
\centering
\caption{Main project problems, solutions, and remaining limitations}
\label{tab:problems-solutions}

\footnotesize
\setlength{\tabcolsep}{4pt}

\begin{tabularx}
{\textwidth}
{>{\raggedright\arraybackslash}p{2.8cm}
 >{\raggedright\arraybackslash}X
 >{\raggedright\arraybackslash}X}
\toprule
Problem
& Action taken
& How far the solution goes \\
\midrule

Nonconsecutive return targets
&
An earlier target construction used the next available observation for each
ticker. This could incorrectly treat a return across a gap of several months as
a one-month return. I added a condition requiring the next observation to be
the exact following calendar month and rebuilt the final panel.
&
This fully solves the calendar-gap problem inside the available monthly
histories. It does not recover returns for months that are completely missing or
delisting returns that are unavailable in Yahoo Finance.
\\[0.3em]

Timing of Welch--Goyal variables
&
An earlier version merged Welch--Goyal variables in the same month as the
stock predictors. I shifted these variables by one month so that stock month
\(t\) uses the macroeconomic state from month \(t-1\).
&
This removes the identified same-month timing problem and makes the rule more
conservative. It does not reproduce the exact historical publication date of
every macroeconomic series.
\\[0.3em]

Timing of accounting information
&
Compustat \texttt{datadate} is the fiscal-year-end date, not the publication
date. I treated each report as available six months after fiscal-year end,
merged it backward to the most recent available report, and prevented very old
reports from being carried forward indefinitely.
&
The procedure substantially reduces look-ahead bias. However, the six-month
delay is still an approximation because exact filing dates are not used.
\\[0.3em]

Missing and extreme predictor values
&
I replaced missing continuous characteristics with the cross-sectional median
of the same month and winsorized these characteristics at the monthly 1st and
99th percentiles. I also applied transparent daily and monthly quality filters
to remove unusable price histories.
&
The final panel contains no missing predictors, infinite values, or duplicate
ticker-month observations. Winsorization limits extreme predictor values, but
the return target remains raw and still contains very large observations.
\\[0.3em]

Inconsistent documentation
&
Earlier documentation described rank normalization, 499 predictors, and an
active Decision Tree. These statements no longer matched the final code. I used
the final repository audit and the authoritative output files to reconcile the
paper with the implemented pipeline.
&
This solves the inconsistency once the paper, README, and Data Construction
Guide are all updated. It does not change the numerical results, but it is
important for making the project understandable and reproducible.
\\[0.3em]


Computational cost of tree tuning
&
Gradient Boosting tuning was very slow because the model fits trees
sequentially across several annual folds. The final specification was chosen
using the separate validation comparison rather than accepting the largest
tuned model.
&
The project reached a usable final specification, but the full tuning process
would be more reproducible if every completed candidate and fold were saved as
a checkpoint.
\\

\bottomrule
\end{tabularx}
\end{table}

The first two issues were particularly important because they could create
look-ahead bias or an incorrect target. After correcting them, I rebuilt the
final data and repeated the later modeling stages. The final panel therefore
uses only consecutive one-month targets and one-month-lagged Welch--Goyal
states.

The accounting solution is less complete. The six-month lag is a practical and
conservative rule, but it does not provide the same precision as exact filing
dates. I therefore treat this as a limitation that has been reduced rather than
fully removed.

\subsection{What worked well}

The strongest methodological decision was the chronological sample design. I
used expanding training windows instead of random cross-validation, compared
fixed and optimized models in a later validation period, and opened the final
test only after the specifications had been chosen.

This structure was useful because it revealed that tuning results from an
earlier period do not necessarily transfer to a later period. The 300-tree
Random Forest and Gradient Boosting candidates performed better within the
training-stage tuning exercise, but worse than their 100-tree versions in the
2015--2019 validation sample. Without the separate validation period, I would
have selected unnecessarily complex final models.

The use of a historical-mean benchmark also worked well. It made the
interpretation direct: a model only provides value under the chosen loss
function if it reduces forecast error relative to a simple constant prediction.

The project also produced useful negative results. Elastic Net selected no
variables, PLS underperformed the benchmark, and Gradient Boosting produced the
worst final test result. These findings are informative because they show that
large predictor sets, regularization, dimension reduction, and nonlinear models
do not automatically generate better forecasts.

\subsection{What the model failures taught me}

Elastic Net shrinks all 484 coefficients to zero. At first, this result looked
like a possible implementation error. After checking the fitted coefficients,
predictions, validation results, and benchmark, I concluded that it is a valid
model outcome. Under the selected penalty, the model preferred a constant
forecast rather than unstable linear effects.

This result does not mean that all stock characteristics are always irrelevant.
It means that the selected Elastic Net specification did not find sufficiently
stable linear information to improve on the historical mean under this sample
and chronology.

PLS presents a different failure. Reducing the 484 predictors to two supervised
components improved validation performance relative to the fixed
20-component version. Nevertheless, it underperformed the benchmark in the
final test. This shows that dimension reduction can control complexity without
guaranteeing that the estimated components remain useful in a later market
period.

Gradient Boosting also illustrates the danger of model flexibility. The model
sequentially corrects earlier errors, but these corrections may fit unstable
patterns. Its negative test out-of-sample \(R^2\) and negative correlation
suggest that the estimated relationships did not transfer successfully to the
test period.

Random Forest is the only complex model with a clear positive full-period
result. Averaging many randomized trees appears to be more effective than
sequential boosting in this sample. However, its success is not stable across
years, so I do not interpret it as a universal relationship.

\subsection{Limits of the Random Forest result}

Random Forest reduces average monthly squared forecast error by approximately
\(1.14\%\), but it has a positive annual out-of-sample \(R^2\) in only one of
the six test years. Its strong performance in 2020 has a large effect on the
full-period ranking.

OLS-3 produces a much smaller total improvement, but it is positive in four of
the six years. This comparison shows that there are at least two ways to assess
a model:

\begin{itemize}
  \item its average performance over the full test period;
  \item the stability of its performance across shorter periods.
\end{itemize}

Random Forest performs better according to the first measure, while OLS-3 is
more frequently above the benchmark. Therefore, the final ranking should not be
interpreted as evidence that Random Forest will always be the best choice.

The target distribution also makes the ranking sensitive to extreme errors.
The raw next-month return contains a small number of very large observations,
and squared loss gives these observations high weight. Reporting MAE provides
a useful comparison, but the primary ranking still depends on MSE.

\subsection{Data limitations that remain}

The stock universe is based on a locked S\&P 500-style ticker list rather than
a fully point-in-time index-constituent database. The locked list ensures that
the same universe is used throughout the pipeline, but it does not fully solve
survivorship and selection bias.

Ticker-based matching can also be affected by:

\begin{itemize}
  \item ticker changes;
  \item mergers and acquisitions;
  \item delistings;
  \item incomplete historical membership information;
  \item differences between Yahoo and Compustat identifiers.
\end{itemize}

Yahoo Finance is convenient and provides a long price history, but it does not
include a permanent security identifier, complete delisting returns, or all the
market microstructure information available from CRSP. Adjusted prices may also
be revised.

The accounting variables are annual, while market variables are monthly or
constructed from daily data. Accounting information therefore changes much
less frequently. Some constructed characteristics are also practical proxies
rather than exact replications of every definition used in the literature.

The 360 characteristic--macro interactions provide flexibility, but they also
create many correlated predictors. Their high Random Forest group importance
partly reflects that this group contains far more variables than the other
groups. It should therefore not be interpreted as proof that every interaction
is economically important.

\subsection{Evaluation limitations}

The test period contains only six years. Although it includes several different
market conditions, it is still a relatively short period for evaluating the
stability of stock-return prediction.

I also do not conduct a formal statistical comparison of forecast losses. The
ranking shows that Random Forest has the lowest monthly MSE, but it does not
show that the difference is statistically significant after accounting for
time dependence and the comparison of several models.

The current out-of-sample \(R^2\) evaluates forecast levels relative to the
historical mean. An additional cross-sectional measure could remove each
month's average return and average prediction. This would test more directly
whether a model can identify which stocks perform relatively better or worse
within each month.

Finally, the project evaluates prediction error and not an investment strategy.
It does not report:

\begin{itemize}
  \item long-short portfolio returns;
  \item Sharpe ratios;
  \item transaction costs;
  \item turnover;
  \item market impact;
  \item short-selling constraints;
  \item factor-adjusted performance.
\end{itemize}

A lower MSE therefore cannot automatically be interpreted as a profitable
trading opportunity.


\subsection{What I would do differently}

With more time and better data access, I would improve the project in the
following order.

I would use a point-in-time CRSP/Compustat database with permanent
identifiers, complete delisting returns, and historical constituent
information. I would also use exact accounting filing dates instead of a fixed
six-month lag.


% ============================================================
% 6. Conclusion
% ============================================================

\section{Conclusion}
\label{sec:conclusion}

This project examined whether machine-learning methods can improve the
prediction of next-month stock excess returns relative to a historical-mean
benchmark. I compared six approaches that represent different levels of model
complexity: the Historical Mean, OLS-3, Partial Least Squares, Elastic Net,
Random Forest, and Gradient Boosting.

The results provide a cautious answer to the research question. Random Forest
ranks first according to the main monthly MSE criterion. Its monthly
out-of-sample \(R^2\) is \(1.1427\%\), meaning that it reduces average monthly
squared prediction error by approximately \(1.14\%\) relative to the historical
mean. This is a positive result, but it is still a modest improvement and should
not be interpreted as an investment return.

OLS-3 ranks second and produces only a very small improvement over the
benchmark. Nevertheless, its result is important because it uses only size,
book-to-market, and 12-month momentum. The fact that such a simple model remains
competitive shows that increasing the number of predictors and the complexity
of the model does not automatically lead to better forecasts.

The other models provide useful negative results. Elastic Net shrinks all 484
coefficients to zero and therefore produces almost the same forecast as the
historical mean. PLS performs worse than the benchmark even after the number of
components is optimized. Gradient Boosting also underperforms and produces a
negative prediction-target correlation. These findings show that
regularization, dimension reduction, and nonlinear modeling are not sufficient
on their own. Their success depends on whether the relationships estimated in
the training period remain stable in later periods.

The yearly analysis also changes the interpretation of the overall ranking.
Random Forest has a positive annual out-of-sample \(R^2\) in only one of the
six test years, while OLS-3 is positive in four years. The Random Forest result
is therefore driven strongly by its performance in 2020. I conclude that
Random Forest has the best average test-period performance, but not that it is
the most reliable model in every market environment.

An important lesson from the project is that the research design matters as
much as the model itself. Expanding-window tuning, a separate validation
period, a simple benchmark, and an untouched final test were necessary to avoid
selecting a model only because it fitted earlier observations well. The
validation results were especially useful because they showed that the more
complex 300-tree Random Forest and Gradient Boosting specifications did not
generalize as well as the fixed 100-tree versions.

The project evaluates statistical forecast accuracy rather than portfolio
performance. The results do not show whether the predictions can generate
profitable long-short portfolios after turnover, transaction costs, liquidity
constraints, and short-selling restrictions. Therefore, the financial
conclusion is limited to forecast performance: Random Forest reduces prediction
error modestly, but the project does not establish a profitable trading
strategy.

A natural extension would be to add neural networks. Neural networks could
capture more complex nonlinear relationships among stock characteristics,
market variables, and macroeconomic states. However, they would also introduce
more hyperparameters and a greater risk of overfitting. I would therefore begin
with a relatively shallow network and use strong regularization, early
stopping, and the same expanding-window validation design used in the current
project.

Because the data contain 484 predictors, I would also examine a
PCA--neural-network approach. Principal Component Analysis could first reduce
the predictor set to a smaller number of components, and the neural network
could then model nonlinear relationships among those components. This may
reduce computational cost, noise, and multicollinearity. At the same time, PCA
is unsupervised and may discard variables that explain little predictor
variance but still contain useful return information. I would therefore compare
three alternatives:

\begin{itemize}
  \item a neural network using the complete predictor set;
  \item a neural network using PCA components;
  \item a neural network using supervised components or selected predictor
        groups.
\end{itemize}

The comparison would have to use the same chronological samples and a genuinely
new test period. I would not select the neural-network architecture based on the
2020--2025 results that have already been observed.

Further extensions would include cross-sectional out-of-sample \(R^2\), formal
tests of differences in prediction loss, and a pre-specified portfolio
evaluation with transaction costs. Better data would also improve the study,
especially a point-in-time CRSP/Compustat universe with permanent identifiers,
delisting returns, and exact accounting filing dates.

Overall, the project does not show that machine learning solves the stock-return
prediction problem. It shows that one ensemble method can produce a modest
improvement, while several other sophisticated methods fail to outperform a
simple benchmark. The main conclusion is therefore that model complexity can
be useful, but only when it is combined with careful timing, chronological
validation, realistic benchmarks, and cautious interpretation.

\FloatBarrier
\printbibliography[heading=bibintoc,title={References}]

% ============================================================
% Appendix
% ============================================================
\clearpage
\appendix



\end{document}
